\documentclass[../../main2.tex]{subfiles}

\begin{document}

\chapter{循环的闭合：预测如何改变预测}
\label{ch:closure}

\begin{motivation}
	上一章遗留的问题：方法给全了，边界也亮了。但预测不是旁观——一份被足够多人相信的预测，会改变被预测的世界本身。我们预测出的未来，会不会反过来塑造新的动机？\\
	\textbf{核心动机}：闭合全书——预测的结果回流为下一轮的动机与目标。并把最后一句话留给诚实：这本书给的是结构，不是预言。
\end{motivation}

\begin{logicmap}
\begin{tikzpicture}[node distance=1.0cm, every node/.style={draw, rectangle, rounded corners, align=center, font=\small}]
\node (q) {预测会改变世界吗？};
\node (s) [below=of q] {自证与自否：两种回音};
\node (b) [below=of s] {回到 Part I：新动机被铸造};
\node (e) [below=of b] {结语：结构，不是预言};
\node (l) [left=2.2cm of b] {Part I\\动机与目标};
\draw[-{Stealth}] (q) -- (s);
\draw[-{Stealth}] (s) -- (b);
\draw[-{Stealth}] (b) -- (e);
\draw[-{Stealth}, dashed] (b) -- (l);
\end{tikzpicture}
\end{logicmap}

\section{自证与自否：预测的两种回音}

\begin{readingnote}
	你有没有想过，为什么央行说「银行很安全」，有时反而引发挤兑？\\
	同一句话，为什么有时救火，有时浇油？
\end{readingnote}

预测一旦公开，就变成了\textbf{新环境}——而环境塑造动机（第 1 章）。于是预测有了两种回音：

\begin{itemize}
	\item \textbf{自证预言}（self-fulfilling prophecy）：预言引发的行为，把预言做成真的。「这支股要涨」传开 $\rightarrow$ 买入 $\rightarrow$ 真的涨。「房价永远涨」信的人多了，抢购把价格顶上去——直到信仰出现裂缝。
	\item \textbf{自否预言}（self-defeating prophecy）：预言引发的行为，把预言作废。「银行要倒」传开 $\rightarrow$ 挤兑 $\rightarrow$ 本来不倒的也倒了——这是自证的反面教材：预言「安全」若被解读为「官方在捂盖子」，同一句话就翻转成自否。
\end{itemize}

\begin{tcolorbox}[colback=green!5, title=\textit{生活类比}]
	天气预报不会改变天气，因为雨不听广播。但「明天菜价要涨」的预报会改变菜价——因为买菜的人听广播。预测的射程，取决于被预测对象\textbf{听不听得见}你。
\end{tcolorbox}

三个回音现场：

\begin{itemize}
	\item \textbf{考试}：老师说「你能考好」 $\rightarrow$ 学生自信、多刷题 $\rightarrow$ 真考好。罗森塔尔实验把这叫期望效应——预言改写了动机。
	\item \textbf{银行挤兑}：1930 年代大萧条的挤兑链、1998 年亚洲金融危机的提款潮——「要倒了」的预言本身是提款指令。
	\item \textbf{流行病}：「板蓝根要断货」的谣言 $\rightarrow$ 抢购 $\rightarrow$ 真断货。预言通过\textbf{模因通道}直接制造了自己的证据。
\end{itemize}

\begin{questionbox}
	\textit{现在你可能想反驳：「按你这么说，预测一公开就作废，那还预测什么？」}\\
	\textbf{不作废，是变复杂。}第 15 章的失败清单里，「反身性」列第三——它加的是难度，不是判死刑。对策是把「预测公开后的行为回流」显式画进链条：预测要连着「预测被谁听见、听见后动不动」一起做。央行学乖了，现在说话字斟句酌——它们在管理\textbf{回音}。
\end{questionbox}

\paragraph*{逻辑衔接}
	回音听清了。现在把全书的环合上：预测不只是被回音骚扰，它\textbf{铸造}新的动机——回到起点。

\section{回到 Part I：预测铸造新动机}

第 1 章说：动机是解释的起点，环境塑造下一轮动机。现在看清了——\textbf{预测也是环境的一部分}。一份被广泛相信的预测（年报、榜单、预言、这本书），会进入千万人的动机清单：

\begin{itemize}
	\item 行业报告说「AI 岗位缺口五百万」$\rightarrow$ 学生转专业、培训班涨价——动机被一份 PDF 改写；
	\item 榜单说「这座城市宜居」$\rightarrow$ 人口流入 $\rightarrow$ 房价涨 $\rightarrow$ 榜单更「准」了——自证又一圈；
	\item \textbf{这本书本身}：它在教你「拆动机、找推手、算贡献」。你信了、用了、讲给朋友听了——社会上就多了一批会拆账的人。预测工具的传播，改变被预测的社会。这个圈，本书自己也逃不掉。
\end{itemize}

所以第 2 章埋的追问（个体加总成什么）和第 7 章的第二枚（加总得到什么），最后的回答是动态的：加总出来的宏观形状，又回流为新的个体动机，再加总……\textbf{这就是「循环」两个字的全部意思}。社会不是一台机器，是一个不停自己改自己的环。

\begin{warningbox}
	\textbf{诚实的终点}：这本书给的是结构，不是预言。它能帮你把一件事「依赖哪几个因素、各占多少」说清楚；它不能告诉你明天会发生什么——尤其在回音最大的地方（人心、市场、舆论），预测的误差被循环放大。能预测的部分，是「不同选择下世界往哪走」；剩下的，留给谦卑。
\end{warningbox}

\paragraph*{逻辑衔接}
	循环闭合，全书收工。最后把武器清单和用途再对一次——结语。

\section{结语：把工具用起来}

六段路，一句话一站：

\begin{itemize}
	\item \textbf{动机}（第 1–2 章）：看人想要什么——目标是排序，行为供出排序。
	\item \textbf{因果}（第 6–7 章）：能预测就等价于有因果——相关是嫌疑，因果是判决。
	\item \textbf{分账}（第 8–9 章）：无你检验起步，交互项单列，份额量级成对。
	\item \textbf{齿轮}（第 10–14 章）：观念、经济、权力、法律——同一动力学的四张脸。
	\item \textbf{循环}（第 15–16 章）：预测改变世界，世界改写动机——用工具时，记得回音。
\end{itemize}

验收还是导论那句话：拿一个你没见过的社会事件，给概率判断，说出依赖哪几个因素。做到了，这本书就完成了它的任务。做不到——回来重看对应那章的「逻辑衔接」，那里标着每一步的接口。

最后一件事，这本书的成功标准其实是第二层的：你看懂了自己。看懂自己为什么焦虑、为什么被算法牵着、为什么「能力强」却没被提拔——共鸣有了，你才会把它转给下一个还在懵的人。裂变的开关，在你看懂自己的那一刻。

\paragraph*{逻辑衔接}
	全书完。数学附录留给想把公式捏在手里的人；致谢留给把这本书从一个念头捏成一叠纸的人。

\begin{thinkbox}
\begin{enumerate}
	\item 挑一个你曾经相信过的「预言」（行业、房价、考试、感情都行）：它是自证了还是自否了？\textbf{你}的行为在回音里占了多大贡献？用无你检验给「如果没有你信它」这个平行宇宙估个数。
	\item \textbf{反事实}：如果这本书真的被一百万人读了，而且他们都学会了「无你检验」——社会会变得更容易预测，还是更难？「拆账的人变多」这件事本身，会怎样回流进你们的动机？把这个环画出来。
	\item \textbf{反事实}：如果所有预测（天气、经济、政策）都禁止公开发布——只许决策者内部使用——世界会更稳还是更乱？谁受益，谁受损？
\end{enumerate}
\end{thinkbox}

\end{document}
